% =============================================================================
% 04_elastic_net.tex — Problem 4: Elastic Net and Neural-Feature Bridge
%                      (P5 — Tài)
% =============================================================================
\section{Elastic Net and Neural-Feature Bridge}
\label{sec:elastic-net}

% ---- 4.1 Research Question and Source Map -----------------------------------
\subsection{Research Question and Source Map}
\label{sec:research-question}

% TODO: Pick ONE research direction to investigate:
%   (a) Weight decay as L2 regularisation in neural networks
%   (b) L1 penalty → sparse weights in neural networks
%   (c) Pruning as post-training sparsification
%   (d) Dropout as implicit regularisation
%
% State your question clearly, e.g.:
% ``How does L2 weight decay in SGD relate to Ridge regression, and what are
%   the differences in the stochastic vs.\ batch setting?''

\textit{TODO: State research question here.}

\begin{table}[H]
\centering
\caption{Source map linking claims to evidence.}
\label{tab:source-map}
\begin{tabularx}{\textwidth}{X l X X}
\toprule
\textbf{Claim} &
\textbf{Source} &
\textbf{Establishes} &
\textbf{Does Not Establish} \\
\midrule
% TODO: Fill in at least 3 rows.
%
% Example row:
% L2 weight decay is equivalent to Ridge regression in batch GD &
% \citet{Krogh1991} &
% Equivalence holds for batch gradient descent with fixed learning rate &
% Does NOT hold for Adam or other adaptive optimisers \\[4pt]
%
\textit{TODO: Claim 1} & \textit{TODO} & \textit{TODO} & \textit{TODO} \\[4pt]
\textit{TODO: Claim 2} & \textit{TODO} & \textit{TODO} & \textit{TODO} \\[4pt]
\textit{TODO: Claim 3} & \textit{TODO} & \textit{TODO} & \textit{TODO} \\
\bottomrule
\end{tabularx}
\end{table}

% ---- 4.2 Elastic Net on Original Predictors ---------------------------------
\subsection{Elastic Net on Original Predictors}
\label{sec:enet-original}

The Elastic Net \citep{ZouHastie2005} combines $\ell_1$ and $\ell_2$ penalties:
\[
    \hat{\bbeta}^{\mathrm{ENet}}
    = \argmin_{\bbeta}
      \left\{
          \frac{1}{2n} \norm{\by - \bX\bbeta}_2^2
          + \lambda\!\left[
              \alpha \norm{\bbeta}_1
              + \frac{1 - \alpha}{2} \norm{\bbeta}_2^2
          \right]
      \right\},
    \quad \alpha \in [0, 1].
\]

% -- Alpha comparison plot
\includefigure{0.75\textwidth}{fig_enet_alpha_comparison}%
    {Elastic Net CV error across different $\alpha$ values.}%
    {fig:enet-alpha}

% -- Elastic Net coefficient plot
\includefigure{0.75\textwidth}{fig_p4_enet_coef}%
    {Elastic Net coefficient estimates at optimal $(\alpha, \lambda)$.}%
    {fig:enet-coef}

% TODO: Report the best (alpha, lambda) combination from CV.
%       Interpret how Elastic Net balances Ridge's grouping effect with
%       Lasso's sparsity.

% ---- 4.3 Fixed ReLU Features -----------------------------------------------
\subsection{Fixed ReLU Features}
\label{sec:relu-features}

To bridge regularised regression with neural-network concepts, we generate
$M = 200$ random features via:
\[
    h_m(\bx) = \max\!\bigl(\bm{w}_m^\top \bx + b_m,\; 0\bigr),
    \qquad m = 1, \dots, M,
\]
where $\bm{w}_m \sim \mathcal{N}(\bm{0}, \bm{I}_p)$ and
$b_m \sim \mathcal{N}(0, 1)$ are \textbf{fixed} (not learned).
Only the output-layer coefficients $\bbeta \in \R^M$ are learned via
Ridge/Lasso/Elastic Net on the transformed design matrix
$\bm{H} \in \R^{n \times M}$.

% -- Neural model CV plot
\includefigure{0.75\textwidth}{fig_neural_cv}%
    {CV error of Ridge/Lasso/Elastic Net on fixed ReLU features ($M = 200$).}%
    {fig:neural-cv}

% TODO: Compare the performance of models trained on original $p = 14$
%       predictors vs.\ the $M = 200$ ReLU features.
%       Discuss: Does the feature expansion help? Why or why not?

% ---- 4.4 Locked Declaration and Final Holdout --------------------------------
\subsection{Locked Declaration and Final Holdout}
\label{sec:holdout}

\bigskip
\noindent
\fbox{\parbox{0.95\textwidth}{%
\textbf{Locked Model Declaration} \\[4pt]
\begin{tabular}{@{}ll@{}}
    \textbf{Feature set:}    & \textit{TODO: original 14 / ReLU 200 / other} \\
    \textbf{Method:}         & \textit{TODO: Ridge / Lasso / Elastic Net} \\
    \textbf{Tuning:}         & \textit{TODO: $\alpha = \ldots$, $\lambda = \ldots$} \\
    \textbf{Justification:}  & \textit{TODO: brief reason} \\
\end{tabular}\\[4pt]
\textit{This declaration was made before examining holdout results.}
}}
\bigskip

% -- Holdout comparison table
\begin{table}[H]
\centering
\caption{Final holdout performance: RMSE and $R^2$ on the test set.}
\label{tab:holdout}
% TODO: Uncomment once generated:
% \input{../output/tables/tab_holdout_comparison}
\fbox{\parbox{0.7\textwidth}{\centering\vspace{0.8cm}
    \texttt{tab\_holdout\_comparison.tex} — run the R pipeline to generate.
\vspace{0.8cm}}}
\end{table}

% -- Actual vs predicted plot
\includefigure{0.65\textwidth}{fig_p4_actual_vs_pred}%
    {Actual vs.\ predicted \texttt{brozek} on the holdout test set.}%
    {fig:actual-vs-predicted}

% TODO: State whether the holdout results support the pre-declared model choice.
%       Discuss any discrepancies.

% ---- 4.5 Deep-Learning Connection and Limitations ---------------------------
\subsection{Deep-Learning Connection and Limitations}
\label{sec:dl-connection}

% TODO: Discuss the following, citing sources from Table~\ref{tab:source-map}:
%
% 1. Weight decay vs.\ L2 penalty:
%    - In batch GD, weight decay ≡ L2 regularisation \citep{Krogh1991}.
%    - In adaptive optimisers (Adam), decoupled weight decay differs
%      \citep{GoodfellowBengioCourville2016}.
%
% 2. L1 penalty → sparse weights:
%    - Lasso encourages sparsity; analogous to L1-regularised neural networks.
%
% 3. Pruning:
%    - Post-training removal of small weights \citep{HanPoolTranDally2015}.
%    - Different from L1: pruning is a discrete, post-hoc operation.
%
% 4. Dropout:
%    - Randomly zeroing activations during training \citep{Srivastava2014}.
%    - Provides an ensemble-like regularisation effect.
%
% Distinguish clearly between these four concepts.

\textit{TODO: Write the deep-learning connection discussion here.}

\paragraph{Limitations.}
% TODO: Discuss at least TWO limitations, e.g.:
%   (a) Fixed ReLU features are not optimised; a trained network would learn
%       better representations.
%   (b) Weight decay equivalence breaks down for adaptive optimisers.
%   (c) L1 penalty in NNs is rarely used in practice due to optimisation
%       difficulties.
%   (d) Our dataset (n=252) is far too small for deep learning to show benefit.

\textit{TODO: Discuss $\geq 2$ limitations.}
